\scanpage{6}
\diagram{\begin{tikzcd}[ampersand replacement=\&,column sep=small,row sep=small]
 \& X\arrow[d,"f"] \&\\
 \& Y=\operatorname{Spec}(A) \&\\
 y\ \text{fermé} \& \& \eta
\end{tikzcd}}
[\uline{strictement}] universellement ouvert, de type fini

$A$ local noethérien \uline{complet}, $f^{-1}(y)$ \inserted{\unsure{ult\textsuperscript{t}|rel\textsuperscript{t}|N.B.}} irréductible

(i) Supposons que \deleted{$f^{-1}(y)$ soit \unsure{ult}} pour $\eta$ générisation convenable de $y$, il y ait \uline{deux} \uline{élts} \uline{distincts} de $E_{X/Y}(\eta)$ qui générisent $x$. En vertu du \uline{lemme} \inserted{a)}, on trouve la $=$ situation, \deleted{avec} par \unsure{changem\textsuperscript{t}} de base $Y'$,
avec \deleted{\unsure{un ann}} $Y'=\operatorname{Spec}(V)$, $V$ anneau de valuation discrète (à corps résiduel alg. clos si on veut, et alors les \deleted{\unsure{fibres}} \inserted{composantes} de $f^{-1}(\eta')$ \uline{ult\textsuperscript{t}} irréductibles), c'est donc qu'il y a \unsure{croisement|croissement|coïncidence} de composantes des fibres géométriques en $f^{-1}(y)$ i.e. en $x$. Supposons que ceci soit exclu.
% slanted
\margin{``\uline{strictement}'' \unsure{univ\textsuperscript{t}} ouvert est utilisé \unsure{seulement|seulem\textsuperscript{t}} dans (i) pour noter que s'il y a \unsure{collapsage|collapse} de fibres \deleted{\unsure{géométriques}} sur $A$, il y en a aussi sur un anneau de valuation discrète --- autrement, on ne s'en sert pas}

(ii) Soit $\xi\in E_{\unsure{\alpha|X}/Y}(\eta)$ tel que $x\in\overline\eta$ [existe, car $f$ est \unsure{univ\textsuperscript{t}} ouvert] je dis que $\overline\xi$ est \uline{ult} irréductible. \deleted{\unsure{On a en}} En effet, on peut supposer $A$ intègre, $\eta$ le p\textsuperscript{t} générique de $Y$. Faisant une extension finie \inserted{séparable} $K'$ sur $K$, \inserted{\unsure{où|rendant} $f^{-1}(\eta)\otimes_K\unsure{K'|K^i}$ décomposé}, et normalisant $A$ dedans, d'où un $A'$ fini sur $A$, on note que $A'$ est \uline{local} (car $A$ local complet). Utilisant le lemme b), \deleted{\unsure{on}} \inserted{et le fait que $f^{-1}(y)$ est \uline{ult\textsuperscript{t} irréductible}, on} trouve qu'on serait dans le cas a), absurde.

On a donc trouvé une application $\varphi:Y\to X$, associant à tout $\eta\in Y$ l'unique \inserted{p\textsuperscript{t}} \deleted{composante irréd.} $\xi$ $\in$ $E_{X/Y}(\unsure{y|\eta})$ tel que $x\in\overline\xi$, et ce $\xi$ est \unsure{ult\textsuperscript{t}|rel\textsuperscript{t}} irréd. Les conditions (i) et (ii) pour une application \unsure{monisante|monisente|monogène} sont satisfaites. Il existe donc un \uline{ouvert} $U$ de $X$ tel que \deleted{pour \unsure{tt}} $U\supset\varphi(Y)$, les fibres de $U$ sur $Y$ \unsure{ét\textsuperscript{t}|étant} ult\textsuperscript{t} irréductibles.
% archive sheet number: 5
